\section{Introduction}

The previous chapter produced training data without a fixed label set, and the introduction to this part explained why a fixed label set is the wrong fit for clinical extraction: the types we want change from one study to the next, and most will never have been seen in training. A model with a fixed output layer cannot produce a type it was not trained on. This chapter builds a model that has no fixed output layer at all. Instead of choosing among a closed list of labels, it compares a span of text against a short description of the type we are looking for and decides whether they match. The type is given at the moment of use, so one model can extract a disease, a dosage, or a treatment response without being retrained for each.

\section{Open-Vocabulary Extraction by Type Description}

This is the idea behind GLiNER~\citep{zaratiana2024gliner}: rather than learning one output per label, the model learns to compare. It reads the document together with the descriptions of the types we are looking for, and scores how well each span of text matches each type. Because the types are supplied as text at inference time, the set of types is open: a user can ask for a type the model never saw in training, and the model will still try to find it.

Our extractor, \emph{ModernCamemBERT-bio-gliner-base}, follows the GLiNER2 design~\citep{zaratiana2025gliner2}. A single encoder reads the document and the candidate type descriptions in one pass, and a light span-scoring head matches every span against every type. The same machinery also handles the other tasks the training data carries: classifying the document, filling structured fields, and extracting relations between spans.

We release the model in two configurations that differ only in their training-data mix. The first is tuned for entity extraction and is the stronger general extractor: it finds types it was never trained on more accurately than any baseline we compare against in \Cref{chap:evaluation}. The second is multi-task: alongside entities it classifies the document, fills structures, and extracts relations, and it does these better than the first, at some cost to entity extraction. The trade-off between the two is itself one of the findings of this part, and we return to it in \Cref{chap:evaluation}.

\section{The MedEmbed Backbone}

The matching is only as good as the encoder that places type descriptions and spans in a shared space. We train our own, \emph{MedEmbed}, a French biomedical sentence encoder. It starts from ModernCamemBERT-bio, the modern French biomedical encoder developed in \Cref{chap:objectives,chap:ontobook}, and is trained with a contrastive objective: a type description and a matching span are pulled together, and mismatched pairs are pushed apart. The model has 149 million parameters and reads up to 8192 tokens, so a long description and a long passage fit in one pass.

The training pairs that matter most come from medical terminologies. The CIM-10, CCAM, and ATC nomenclatures are turned into description-span pairs by the ontology work of \Cref{chap:ontobook}, which we do not repeat here.

On OntoBench-FR, a retrieval benchmark built from the French terminologies and from real clinical cases, MedEmbed reaches 58.1\% average accuracy and beats encoders several times its size: 48.9\% for mE5-large-instruct, 47.2\% for BGE-M3, and 44.9\% for Solon.\footnote{This is the average weighted by the number of test items per benchmark; an unweighted average gives a slightly higher figure on a different scale, and the two should not be compared directly.} An ablation confirms where the strength comes from: the terminology pairs are the one source that matters, and training on them alone is enough to match the full mixture, while removing them costs more than ten points. The starting encoder matters too: ModernCamemBERT-bio gives a clear gain over a plainly pretrained encoder on the novel types that generalisation depends on.

\section{Training}

We train the extractor on the synthetic annotations of \Cref{chap:synthetic}, with no real clinical labels at any point. Both configurations use the same recipe and run in a few hours on a single GPU. The result is a small open model that extracts types defined at inference time, trained entirely on public and synthetic data.

\section*{Conclusion}

A model that can extract a type it was never trained on raises a problem the moment we try to measure it. The usual way to score extraction is to compare against a fixed human annotation, but a fixed annotation only covers the types the annotators were asked to mark, which is exactly the closed world this model was built to escape. Measuring it fairly needs a different kind of evaluation, and that is the subject of the next chapter.
